\documentclass{article}
\usepackage{spconf,amsmath,graphicx}
\usepackage{cite}
\usepackage{amssymb,amsfonts}
\usepackage{algorithmic}
\usepackage{longtable}
\usepackage{nth}
\usepackage{multirow}
\usepackage{multicol}
\usepackage{rotating}
\usepackage{textcomp}
\usepackage[dvipsnames]{xcolor}
\usepackage[vietnamese,main=english]{babel}
\usepackage{enumitem}
\usepackage{float}
\usepackage{hyperref} % hyperref must be loaded last
\def\BibTeX{{\rm B\kern-.05em{\sc i\kern-.025em b}\kern-.08em
    T\kern-.1667em\lower.7ex\hbox{E}\kern-.125emX}}

% Force LaTeX to place floats near their source text
\renewcommand{\topfraction}{0.95}
\renewcommand{\bottomfraction}{0.95}
\renewcommand{\textfraction}{0.05}
\renewcommand{\floatpagefraction}{0.9}
\setcounter{topnumber}{4}
\setcounter{bottomnumber}{3}
\setcounter{totalnumber}{6}
\setcounter{dbltopnumber}{3}
\renewcommand{\dbltopfraction}{0.95}

\begin{document}
\ninept

\title{A Hard Negative-Aware Optimization for Multilingual Text-Based Person Search}

\name{Tung-Lam Pham$^{\star}$, Thi-Hoai Phan$^{\star}$, Thuy-Binh Nguyen$^{\dagger}$, Thanh-Hai Tran$^{\star}$, \\ Ngoc-Diep Do$^{\star}$, Hong-Quan Nguyen$^{\ddagger}$, Thi-Lan Le$^{\star}$}
\address{$^{\star}$~SEEE, Hanoi University of Science and Technology, Vietnam \\
$^{\dagger}$~Faculty of Electrical-Electronic Eng., University of Transport and Communications, Vietnam \\
$^{\ddagger}$~Faculty of Information Technology, Viet-Hung Industrial University, Vietnam}

\maketitle

\begin{abstract}
Multilingual Text-Based Person Search (TBPS) remains challenging in low-resource settings due to weak cross-modal alignment. While TBPS-mSigLIP improves robustness with sigmoid-based contrastive learning, it lacks effective discrimination for hard negatives. We propose mSigLIP-CLoRA, a parameter-efficient framework that enhances TBPS-mSigLIP by integrating Cross-modal Circle Loss for fine-grained discrimination and LoRA for stable, efficient training. Experiments show that our method achieves state-of-the-art performance in data-scarce scenarios, reaching 52.28\% Rank@1 on 3000VnPersonSearch and significantly outperforming baselines with limited data, while remaining competitive on large-scale benchmarks with far fewer trainable parameters. The source code is available at \href{https://github.com/pahmlam/Research_on_CircleLoss_for_TBPS-mSigLIP.git}{our code repository}.
\end{abstract}


\begin{keywords}
Text-Based Person Search, Multilingual Search, Circle Loss, Low-Rank Adaptation, mSigLIP
\end{keywords}

\section{Introduction}\label{intro}

Multilingual Text-Based Person Search (MTBPS) aims to retrieve pedestrian images using natural language descriptions, which is particularly valuable when visual queries are unavailable, such as in surveillance and forensic scenarios \cite{Li_2017_CVPR}. While recent TBPS methods have achieved strong performance on English benchmarks, extending them to multilingual and low-resource settings remains challenging due to limited annotations and weak cross-lingual alignment \cite{pmlr-v139-jia21b,9648695}.

Recent Vision-Language Models (VLMs), especially CLIP-based architectures and multilingual variants such as mSigLIP \cite{Zhai_2023_ICCV}, have significantly advanced TBPS by learning shared embedding spaces via contrastive learning. TBPS-mSigLIP \cite{tran2025multilingual} further adopts a Sigmoid-based Noise-robust Image-Text Contrastive (N-ITC) loss to mitigate label noise in multilingual training. However, we observe that Sigmoid-based objectives provide insufficient gradient differentiation between easy and hard negatives, resulting in weak intra-class compactness and suboptimal inter-class separation—an issue that becomes more pronounced in low-resource regimes.

To address this limitation, we propose \textbf{mSigLIP-CLoRA}, a hard-negative-aware and parameter-efficient optimization framework built upon the TBPS-mSigLIP baseline. Specifically, we introduce an auxiliary \textit{Cross-modal Circle Loss} \cite{sun2020circle} to enhance fine-grained discrimination through adaptive pair-wise re-weighting. To balance global alignment and local geometric refinement, we investigate both fixed and curriculum-based weighting strategies. Moreover, to stabilize training under explicit hard-negative optimization, we incorporate \textit{Low-Rank Adaptation (LoRA)} \cite{hu2022lora}, which constrains updates to a low-rank subspace, improving robustness and enabling larger batch sizes with significantly fewer trainable parameters.

Our contributions: (1) We analyze the optimization limitations of Sigmoid-based contrastive losses for hard-negative modeling in multilingual TBPS. (2) We propose a unified framework integrating Cross-modal Circle Loss with LoRA, with a systematic study of weighting strategies. (3) Extensive experiments demonstrate state-of-the-art results on 3000VnPersonSearch and significant gains on the 10\% CUHK-PEDES subset, with competitive performance on high-resource benchmarks using superior parameter efficiency.


\section{Related Work}\label{related}

Text-Based Person Search (TBPS) has evolved from early dual-encoder CNN–RNN frameworks optimized with ranking-based losses \cite{Zhang_2018_ECCV,Li_2017_CVPR} to modern vision--language models trained via large-scale contrastive learning \cite{pmlr-v139-radford21a}. CLIP-style approaches demonstrate strong generalization, while recent studies further enhance cross-modal alignment through fine-grained matching and implicit reasoning mechanisms \cite{jiang2023cross,yin2025cross}. However, most existing TBPS methods primarily focus on English benchmarks. To address multilingual and low-resource scenarios, TBPS-mSigLIP \cite{tran2025multilingual} introduces a Sigmoid-based Noise-robust Image--Text Contrastive (N-ITC) loss, improving robustness to label noise and memory efficiency. Despite these advantages, its optimization emphasizes global alignment and lacks explicit mechanisms for hard-negative discrimination, which limits the separability of the learned embedding space.

The choice of metric learning objective plays a critical role in retrieval performance \cite{kulis2013metric}. While Sigmoid-based contrastive losses scale well, they provide limited gradient separation between easy and hard negatives. Traditional ranking losses, such as Triplet Loss \cite{schroff2015facenet}, rely on hard mining but often suffer from unstable optimization. Angular margin-based objectives, notably Circle Loss \cite{sun2020circle}, address these issues via adaptive pair-wise re-weighting, explicitly emphasizing hard samples and refining embedding geometry. Although Circle Loss has proven effective in unimodal retrieval, its application to multilingual cross-modal person search remains largely unexplored.

In parallel, Parameter-Efficient Fine-Tuning (PEFT) techniques have emerged to reduce overfitting and computational cost in low-data regimes. Low-Rank Adaptation (LoRA) \cite{hu2022lora} freezes backbone parameters while introducing low-rank updates, significantly reducing trainable parameters and improving training stability. In this work, LoRA complements Circle Loss by regularizing hard-negative optimization and enabling larger batch sizes, which are essential for effective contrastive learning.

\section{Proposed Method}\label{proposed}
\begin{figure*}[!t]
    \centering
    \includegraphics[width=0.85\linewidth]{Figure/framework.png} 
    \caption{Overall architecture of the proposed \textbf{mSigLIP-CLoRA} framework for Multilingual Text-Based Person Search.}
    \label{fig:framework}
\end{figure*}

This section presents \textbf{mSigLIP-CLoRA}, a hard-negative-aware and parameter-efficient optimization framework built upon TBPS-mSigLIP \cite{tran2025multilingual} . While TBPS-mSigLIP provides robust global alignment through a Sigmoid-based contrastive objective, it lacks explicit geometric constraints for separating hard negatives. To address this limitation, we integrate an auxiliary \textit{Cross-modal Circle Loss} for fine-grained geometric refinement and adopt \textit{Low-Rank Adaptation (LoRA)} to stabilize optimization and improve training efficiency. The overall architecture is illustrated in Fig.~\ref{fig:framework}.

\subsection{Baseline: TBPS-mSigLIP Framework}
Our framework is constructed upon the strong baseline of TBPS-mSigLIP \cite{tran2025multilingual}, a state-of-the-art approach for multilingual person search. The core model adopts a Vision Transformer (ViT-B/16) as the image encoder and a multilingual Transformer as the text encoder. Unlike CLIP, which relies on a Softmax-based InfoNCE loss, mSigLIP formulates image--text alignment using a pairwise sigmoid loss, treating matching as independent binary classification tasks. This design significantly reduces memory consumption and improves training stability under large batch sizes. To adapt the foundation model to the person search task, the baseline optimizes a multi-task objective:
\begin{equation}
\mathcal{L}_{base}
= \alpha_1 \mathcal{L}_{N\text{-}ITC}
+ \alpha_2 \mathcal{L}_{MVS}
+ \alpha_3 \mathcal{L}_{C\text{-}ITC}
+ \alpha_4 \mathcal{L}_{SS}.
\end{equation}

\noindent\textbf{Notation.}
Let $N$ denote the mini-batch size; $\mathbf{v}_i, \mathbf{u}_i \in \mathbb{R}^d$ are the $L_2$-normalized image and text embeddings, respectively.
$\text{sim}(\cdot,\cdot)$ denotes cosine similarity, $z_{ij}\!\in\!\{+1,-1\}$ indicates matched pairs, and $\gamma, c$ are the scale and bias parameters of the sigmoid formulation.

\begin{itemize}[nosep,leftmargin=*]
\item $\mathcal{L}_{N\text{-}ITC}$ (Normalized Image--Text Contrastive Loss) aligns global image and text representations while mitigating noisy annotations:
\vspace{-0.5em}
\begin{equation}
\mathcal{L}_{N\text{-}ITC}= -\frac{1}{N}\sum_{i=1}^{N} \sum_{j=1}^{N}\log\sigma \!\left(z_{ij} \left( \gamma \, \mathbf{v}_i^\top \mathbf{u}_j - c \right)\right)
\end{equation}

\vspace{-0.5em}
\item $\mathcal{L}_{C\text{-}ITC}$ (Cyclic Image--Text Contrastive Loss) enforces structured embedding alignment via in-modality consistency and cross-modality symmetry:
\vspace{-0.5em}
\begin{equation}
\mathcal{L}_{C^I\text{-}ITC}\!=\!\frac{1}{N^2}\!\sum_{i,j}\!\big(\text{sim}(\mathbf{v}_i, \mathbf{v}_j)\!-\!\text{sim}(\mathbf{u}_i, \mathbf{u}_j)\big)^2
\end{equation}
\vspace{-1em}
\begin{equation}
\mathcal{L}_{C^C\text{-}ITC}\!=\!\frac{1}{N^2}\!\sum_{i,j}\!\big(\text{sim}(\mathbf{v}_i, \mathbf{u}_j)\!-\!\text{sim}(\mathbf{v}_j, \mathbf{u}_i)\big)^2
\end{equation}
\vspace{-1em}
\begin{equation}
\mathcal{L}_{C\text{-}ITC} = \lambda_I \mathcal{L}_{C^I\text{-}ITC} + \lambda_C \mathcal{L}_{C^C\text{-}ITC},
\end{equation}
with $\lambda_I = \lambda_C = 0.25$.

\vspace{-0.3em}
\item $\mathcal{L}_{MVS}$ (Multi-View Supervision Loss) enforces consistency between original images and their augmented views using the N-ITC formulation.

\vspace{-0.3em}
\item $\mathcal{L}_{SS}$ (Self-Supervision Loss) promotes representation invariance between two augmented views of the same image:
\vspace{-0.5em}
\begin{equation}
\mathcal{L}_{SS}=-\frac{1}{2N}\sum_{i=1}^{2N}\log\frac{\exp \!\left(\text{sim}(\mathbf{v}_i, \mathbf{v}_{i^+})/\tau_s\right)
}{\sum_{k \neq i}\exp \!\left(\text{sim}(\mathbf{v}_i, \mathbf{v}_k)/\tau_s\right)}
\end{equation}
\end{itemize}

While this composite objective provides strong supervision, the primary alignment loss $\mathcal{L}_{N\text{-}ITC}$ yields relatively uniform gradients, offering limited control over the relative importance of hard and easy negative pairs during fine-grained optimization.

\subsection{Auxiliary Cross-Modal Circle Loss}
To enhance fine-grained discrimination without compromising the global alignment of the foundation model, we propose an \textbf{Auxiliary Cross-Modal Circle Loss} ($\mathcal{L}_{\text{circle}}$) operating in parallel with the primary N-ITC loss. Let $\mathcal{P}$ and $\mathcal{N}$ denote the sets of positive and negative pairs, and $s$ be the cosine similarity. We formulate the loss to maximize within-class ($s_p$) and minimize between-class ($s_n$) similarity:
\begin{equation}
\mathcal{L}_{\text{circle}} = \log \Bigg[ 1 + \sum_{j \in \mathcal{N}} e^{\gamma \alpha_n^j (s_n^j - m)} \cdot \sum_{i \in \mathcal{P}} e^{-\gamma \alpha_p^i (s_p^i - (1-m))} \Bigg],
\end{equation}
where $\gamma$ is the scale factor and $m$ is the relaxation margin. Crucially, we employ adaptive weights $\alpha_p^i = [1 + m - s_p^i]_+$ and $\alpha_n^j = [s_n^j + m]_+$ ($[\cdot]_+$ is ReLU) to dynamically amplify gradients for hard samples while suppressing well-separated pairs \cite{sun2020circle}.

\textbf{Optimization Objective and Scheduling:} The final objective is $\mathcal{L} = \mathcal{L}_{\text{base}} + \alpha_5(t) \cdot \mathcal{L}_{\text{circle}}$, where $\mathcal{L}_{\text{base}}$ aggregates N-ITC, C-ITC, MVS, and SS losses \cite{tran2025multilingual}. We investigate two strategies for $\alpha_5(t)$: (1) \textbf{Static Weighting}, where $\alpha_5$ is constant; and (2) \textbf{Curriculum Learning}, designed to prevent early disruption of global alignment. In the latter, $\alpha_5(t)$ remains 0 for $T_{warmup}$ epochs, then linearly increases to $\alpha_{max}$ over a period of $T_{ramp}$, balancing stability and discrimination.

\subsection{Gradient Stabilization via LoRA}
Angular-margin objectives introduce high-magnitude gradients that can destabilize fine-tuning. To address this, we apply Low-Rank Adaptation (LoRA) to the query, key, and value projections while freezing the backbone. LoRA constrains updates to a low-rank subspace, stabilizing optimization by buffering sharp gradients, and significantly reduces memory usage, enabling a $3\times$ larger batch size. The increased batch implicitly enlarges the negative set, which is crucial for effective hard-negative penalization in Circle Loss. 


\section{Experiments and Results}\label{eexp}

\subsection{Datasets}
We evaluate on three TBPS benchmarks spanning three typologically diverse languages: CUHK-PEDES \cite{Li_2017_CVPR} (English, high-resource), 3000VnPersonSearch \cite{pham2022towards} (Vietnamese, low-resource), and PRW-TPS-CN \cite{wu2023prototype} (Chinese, logographic script). \textbf{CUHK-PEDES} contains 40,206 images of 13,000 identities with 80,440 English descriptions; 34,054/3,074 images for training/testing. \textbf{3000VnPersonSearch} has 6,302 images of 3,000 identities with two Vietnamese descriptions each; 4,302/2,000 for training/testing. \textbf{PRW-TPS-CN} is built upon the PRW person search dataset with 47,102 Chinese text descriptions, enabling cross-script evaluation on a logographic writing system.

\subsection{Evaluation Metrics}
We adopt Rank-$k$ accuracy, mean Average Precision (mAP), and mean Inverse Negative Penalty (mINP) following standard TBPS evaluation protocols.

\subsection{Implementation Details}
All experiments use PyTorch on the TBPS-mSigLIP framework with pre-trained mSigLIP weights. Images are resized to $256 \times 256$. The baseline \cite{tran2025multilingual} uses full fine-tuning of all 376M parameters with batch size 8 on a single NVIDIA RTX 3060 GPU. mSigLIP-CLoRA applies LoRA with rank $r\!=\!32$, scaling $\alpha\!=\!64$ (effective ratio $\alpha/r\!=\!2.0$), and dropout $0.05$ to the query, key, value, and output projections of all transformer layers, yielding 5.9M trainable parameters (1.57\%). Training runs for 60 epochs with AdamW at batch size 24. Loss weights: $\alpha_1\!=\!1.0$ (N-ITC), $\alpha_2\!=\!1.0$ (MVS), $\alpha_3\!=\!0.1$ (C-ITC), $\alpha_4\!=\!0.4$ (SS) \cite{tran2025multilingual}, plus auxiliary Circle Loss weight $\alpha_5$ under fixed and curriculum schedules (Section~\ref{sec:curriculum_ablation}).

\subsection{Ablation Study}

\subsubsection{Impact of LoRA and Circle Loss}
We analyze the individual and combined effects of LoRA and Circle Loss on the 3000VnPersonSearch dataset, with results summarized in Table~\ref{tab:ablation_main}.

\begin{table}[H]
    \centering
    \vspace{-0.5em}
    \caption{Ablation study on the impact of LoRA and Circle Loss on the 3000VnPersonSearch dataset. Best results are in \textbf{bold}.}
    \label{tab:ablation_main}
    \resizebox{\columnwidth}{!}{%
    \begin{tabular}{|l|c|c|c|c|c|c|}
        \hline
        \textbf{Methods} & \textbf{Batch} & \textbf{Circle} & \textbf{R@1} & \textbf{R@5} & \textbf{R@10} & \textbf{mINP} \\ \hline
        TBPS-mSigLIP (Full FT) \cite{tran2025multilingual} & 8 & - & 49.70 & 75.93 & 84.75 & 48.66 \\ \hline
        LoRA Only & 24 & - & 49.90 & \textbf{78.05} & 86.30 & \textbf{49.45} \\ \hline
        LoRA + Circle (Fixed $\alpha_5=0.1$) & 24 & \checkmark & \textbf{50.53} & 77.78 & \textbf{86.43} & 49.37 \\ \hline
    \end{tabular}%
    }
\end{table}

Replacing full fine-tuning with LoRA enables a $3\times$ batch size increase, improving contrastive learning. Adding Circle Loss further boosts Rank@1 to 50.53\%, confirming the complementary roles of LoRA and Circle Loss.

\subsubsection{Isolating LoRA Regularization from Batch-Size Effects}
To disentangle LoRA's regularization benefit from the batch-size increase, we conduct a control experiment: full fine-tuning of all 376M parameters at the identical batch configuration as LoRA (batch=24, gradient accumulation=3, effective batch=72) on an RTX 4090. Results on 3000VnPersonSearch are shown in Table~\ref{tab:control_ft}.

\begin{table}[H]
    \centering
    \vspace{-0.5em}
    \caption{Control experiment isolating LoRA regularization from batch-size effects on 3000VnPersonSearch. All methods use effective batch size 72.}
    \label{tab:control_ft}
    \resizebox{\columnwidth}{!}{%
    \begin{tabular}{|l|c|c|c|c|c|c|}
        \hline
        \textbf{Method} & \textbf{Params} & \textbf{R@1} & \textbf{R@5} & \textbf{R@10} & \textbf{mAP} & \textbf{mINP} \\ \hline
        Full FT (no Circle) & 376M (100\%) & 49.18 & 76.30 & 85.58 & 54.49 & 47.87 \\ \hline
        LoRA (no Circle) & 5.9M (1.57\%) & 49.90 & 77.45 & 86.20 & 55.23 & 48.65 \\ \hline
        \textbf{LoRA + Curriculum Circle} & 5.9M (1.57\%) & \textbf{52.28} & \textbf{79.55} & \textbf{88.03} & \textbf{57.32} & \textbf{50.57} \\ \hline
    \end{tabular}%
    }
    \vspace{-1em}
\end{table}

At the identical batch configuration, full fine-tuning (49.18\%) underperforms LoRA (49.90\%), ruling out batch size as the source of improvement. With 376M trainable parameters, full fine-tuning overfits on the small VN3K dataset. The improvement from Circle Loss is +2.38\% R@1 (49.90\% $\to$ 52.28\%), cleanly isolating its contribution. This three-way comparison confirms that both LoRA regularization and Circle Loss hard-negative mining contribute independently.

\subsubsection{Sensitivity Analysis of Circle Loss Weight}
To determine the optimal contribution of the geometric refinement, we conducted a sensitivity analysis on the Circle Loss weight ($\alpha_5$) using a fixed strategy. Table \ref{tab:circle_weight} summarizes the performance across different weight settings.

\begin{table}[H]
    \centering
    \vspace{-0.5em}
    \caption{Sensitivity analysis of Circle Loss weight ($\alpha_5$) on 3000VnPersonSearch with LoRA enabled (Batch Size = 24).}
    \label{tab:circle_weight}
    \resizebox{\columnwidth}{!}{%
    \begin{tabular}{|c|c|c|c|c|c|} 
        \hline
        $\alpha_5$ & \textbf{R@1} & \textbf{R@5} & \textbf{R@10} & \textbf{mAP} & \textbf{mINP} \\ \hline
        0.0 (LoRA only) & 49.90 & \textbf{78.05} & 86.30& 55.83 & 49.45 \\ \hline
        0.1 & \textbf{50.53} & 77.78 & 86.43 & \textbf{55.94} & 49.37 \\ \hline
        0.2 & 49.83 & 77.33 & 86.43 & 55.46 & 49.02 \\ \hline
        0.3 & 49.83 & 77.23 & \textbf{86.73} & 55.68 & \textbf{49.45} \\ \hline
    \end{tabular}%
    }
\end{table}

The best performance is achieved at $\alpha_5=0.1$. For larger values ($\alpha_5 \ge 0.2$), performance degrades. We attribute this instability to two primary factors: (1) \textbf{High gradient variance}: With a relatively small batch size of 24, the "hardest negative" fluctuates significantly across mini-batches, creating unstable optimization targets. (2) \textbf{Gradient amplification}: The scaling factor $\gamma=128$ produces exceptionally large gradients that, when weighted heavily, can exceed the limited regularization capacity of the LoRA subspace and destabilize the fine-tuning process.

\subsubsection{Effectiveness of Curriculum Learning} \label{sec:curriculum_ablation}
Building upon the sensitivity analysis, we hypothesize that applying Circle Loss with a fixed weight from initialization may disrupt the early alignment of LoRA parameters. To mitigate this, we propose a \textbf{Curriculum Hard-Mining Schedule}, where $\alpha_5$ linearly increases from 0 to 0.1 over epochs 6–20 (after a 5-epoch warm-up).

\begin{table}[H]
    \centering
    \vspace{-0.5em}
    \caption{Impact of the Progressive Hard-Mining Schedule on 3000VnPersonSearch.}
    \label{tab:curriculum}
    \resizebox{\columnwidth}{!}{%
    \renewcommand{\arraystretch}{1.2}
    \begin{tabular}{|l|c|c|c|c|c|}
        \hline
        \textbf{Strategies} & \textbf{R@1} & \textbf{R@5} & \textbf{R@10} & \textbf{mAP} & \textbf{mINP} \\ \hline
        Fixed Weight ($\alpha_5=0.1$) & 50.53 & 77.78 & 86.43 & 55.94 & 49.37 \\ \hline
        \textbf{Curriculum Schedule} & \textbf{52.28} & \textbf{79.55} & \textbf{88.03} & \textbf{57.32} & \textbf{50.57} \\ \hline
        \textit{Improvement} & \textit{+1.75\%} & \textit{+1.77\%} & \textit{+1.60\%} & \textit{+1.38\%} & \textit{+1.20\%} \\ \hline
    \end{tabular}%
    \renewcommand{\arraystretch}{1}
    }
    \vspace{-1.5em}
\end{table}

As shown in Table \ref{tab:curriculum}, this strategy achieves a new state-of-the-art Rank-1 of \textbf{52.28\%}, confirming that a warm-up period for global alignment allows stable semantic structure formation before Circle Loss constraints. Specifically, LoRA parameters are initialized near zero and require time to establish alignment. The 5-epoch warm-up ensures the primary N-ITC loss converges to a reasonable structure, while the subsequent linear ramp over 15 epochs provides a predictable integration that avoids sudden gradient shocks, preventing premature hard-negative mining from destabilizing the early-stage representations.
\begin{figure*}[!t]
    \centering
    \includegraphics[width=0.95\linewidth]{Figure/flipped_cases_visualization.png}   
    \caption{Qualitative comparison on 3000VNPersonSearch. \textcolor{ForestGreen}{Green} = correct matches; \textcolor{red}{red} = incorrect retrievals.}
    \label{fig:qualitative}   
\end{figure*}

We visualize the similarity distribution in Fig.~\ref{fig:sim_distribution} by plotting positive ($s_p$) vs.\ hardest negative ($s_n$) cosine similarities.

\begin{figure}[H]
    \centering
    \includegraphics[width=\columnwidth]{Figure/distribution_final_v5_pub.png} 
    \caption{\textbf{Geometric Analysis of Similarity Distribution} ($s_n$ vs. $s_p$). (a) Baseline TBPS-mSigLIP (b) Our method (LoRA + Curriculum Circle Loss).}
    \label{fig:sim_distribution}
\end{figure}


To explain the observed performance gains, we analyze the gradient behavior of N-ITC and Circle Loss with respect to similarity scores. Fig.~\ref{fig:gradients} presents an analytical comparison that is independent of data distributions.

\begin{figure}[H]
    \centering
    \includegraphics[width=\linewidth]{Figure/gradient_3d_optimized_pub.png}
    \caption{\textbf{Theoretical visualization of gradient magnitude.}
    (a) Gradient w.r.t.\ negative similarity ($s_n$).
    (b) Gradient w.r.t.\ positive similarity ($s_p$).}
    \label{fig:gradients}
\end{figure}


For negative pairs, N-ITC provides weak signals once similarities fall below moderate levels. Circle Loss introduces margin-aware gradients, ensuring hard negatives receive strong penalization. For positive pairs, N-ITC gradients diminish rapidly as similarity increases, while Circle Loss sustains non-negligible gradients even for highly similar pairs, continuously tightening intra-class clusters.

As shown in Fig.~\ref{fig:sim_distribution}(a), the baseline exhibits near-linear convergence toward the decision boundary ($y\!=\!x$). In contrast, Fig.~\ref{fig:sim_distribution}(b) shows our method enforces a spherical geometry (\textcolor{blue}{blue} curve) aligning with the theoretical margin ($m\!=\!0.25$), suppressing hard negative overlaps. The 52.6\% Rank-1 corresponds to FP32 inference, slightly higher than the FP16 result in Table~\ref{tab:sota_vn}.
\subsection{Comparison with State-of-the-Art Methods}
\subsubsection{Results on Low-Resource Dataset (3000VnPersonSearch)}
Table \ref{tab:sota_vn} presents the performance comparison on the Vietnamese dataset. Our mSigLIP-CLoRA (with Curriculum Learning) outperforms existing approaches by a significant margin.

\begin{table}[H]
    \centering
    \vspace{-0.5em} 
    \caption{Comparison with SOTA methods on 3000VnPersonSearch. Best results are in bold.}
    \label{tab:sota_vn}
    \resizebox{0.9\columnwidth}{!}{%
    \begin{tabular}{|l|c|c|c|c|c|}
        \hline
        \textbf{Methods} & \textbf{R@1} & \textbf{R@5} & \textbf{R@10} & \textbf{mAP} & \textbf{mINP} \\ \hline
        SRCF \cite{phan2023method} & 30.70 & 57.05 & 67.93 & - & - \\ \hline
        CFFVL \cite{phan2023method} & 40.32 & 66.60 & 76.58 & - & - \\ \hline
        ViTAA \cite{phan2024systematic} & 27.08 & 51.38 & 63.00 & - & - \\ \hline
        TBPS-mSigLIP \cite{tran2025multilingual} & 49.70 & 75.93 & 84.75 & 54.96 & 48.66 \\ \hline
        \textbf{mSigLIP-CLoRA} & \textbf{52.28} & \textbf{79.55} & \textbf{88.03} & \textbf{57.32} & \textbf{50.57} \\ \hline
    \end{tabular}%
    }
\end{table}

To verify result stability, we re-run the main experiment (LoRA + Curriculum Circle Loss) with 3 different random seeds: R@1 = 51.52 $\pm$ 0.68\%, R@5 = 78.78 $\pm$ 0.69\%, R@10 = 87.22 $\pm$ 0.71\%, mAP = 56.95 $\pm$ 0.44\%, mINP = 50.56 $\pm$ 0.67\%. The low standard deviation ($<$0.75\%) across all metrics confirms that improvements are stable and not attributable to random seed selection.

\subsubsection{Results on High-Resource Dataset (CUHK-PEDES)}
To verify generalizability, we evaluate on the large-scale CUHK-PEDES dataset. Table \ref{tab:sota_cuhk} compares our results with SOTA methods.

\begin{table}[H]
    \centering
    \caption{Comparison on the CUHK-PEDES dataset. Best results appear in bold; second-best results are underlined.}
    \label{tab:sota_cuhk}
    \resizebox{\columnwidth}{!}{%
    \begin{tabular}{|l|c|c|c|c|c|}
        \hline
        \textbf{Methods} & \textbf{R@1} & \textbf{R@5} & \textbf{R@10} & \textbf{mAP} & \textbf{mINP} \\ \hline
        SAF \cite{li2022learning} & 64.13 & 82.62 & 88.40 & - & - \\ \hline
        IRRA \cite{jiang2023cross} & \underline{73.38} & \textbf{89.93} & \textbf{93.71} & \textbf{66.13} & \textbf{50.24} \\ \hline
        TBPS-CLIP \cite{cao2024empirical} & \textbf{73.54} & 88.19 & 92.35 & \underline{65.38} & - \\ \hline
        TBPS-mSigLIP \cite{tran2025multilingual} & 70.76 & 87.44 & 92.06 & 63.40 & 46.92 \\ \hline
        \textbf{mSigLIP-CLoRA} & 71.85 & \underline{88.26} & \underline{92.98} & 64.34 & \underline{48.11} \\ \hline
    \end{tabular}%
    }
\end{table}

While our method improves upon the baseline across all metrics (e.g., Rank-1 increases from 70.76\% to 71.85\%), the gain (+1.09\%) is relatively modest compared to the low-resource Vietnamese setting (+2.58\%). We attribute this to the \textit{false negative collision rate}. Large-scale datasets like CUHK-PEDES contain many distinct individuals wearing similar clothing in similar environments. Under aggressive hard-negative mining with a static margin ($m=0.25$), these visually identical but distinct identities are frequently sampled as the "hardest negative." Forcing a large margin between them leads to over-penalization, whereas low-resource datasets have fewer such collisions.

\subsubsection{Results on Chinese Dataset (PRW-TPS-CN)}
To validate cross-script generalization, we evaluate on PRW-TPS-CN with Chinese captions. Table~\ref{tab:sota_cn} confirms strong improvements across a logographic writing system.

\begin{table}[H]
    \centering
    \caption{Comparison on PRW-TPS-CN (Chinese). Best results in bold.}
    \label{tab:sota_cn}
    \resizebox{\columnwidth}{!}{%
    \begin{tabular}{|l|c|c|c|c|c|}
        \hline
        \textbf{Methods} & \textbf{R@1} & \textbf{R@5} & \textbf{R@10} & \textbf{mAP} & \textbf{mINP} \\ \hline
        TPAN \cite{wu2023prototype} & 21.63 & 42.54 & 52.99 & - & - \\ \hline
        TBPS-mSigLIP \cite{tran2025multilingual} & 46.78 & 60.28 & 66.82 & 35.41 & 10.61 \\ \hline
        \textbf{mSigLIP-CLoRA} & \textbf{59.35} & \textbf{70.58} & \textbf{75.48} & \textbf{46.44} & \textbf{15.10} \\ \hline
    \end{tabular}%
    }
\end{table}

\subsubsection{Generalizability on Low-Data Regimes (10\% CUHK-PEDES)}
To evaluate robustness under data scarcity, we conduct experiments on a 10\% subset of CUHK-PEDES. Table~\ref{tab:cuhk_10percent} compares the baseline with Fixed Weight and Curriculum strategies.

\begin{table}[H]
    \centering
    \caption{Performance comparison on 10\% of the CUHK-PEDES training set. }
    \label{tab:cuhk_10percent}
    \resizebox{\columnwidth}{!}{%
    \begin{tabular}{|l|c|c|c|c|c|}
        \hline
        \textbf{Method} & \textbf{R@1} & \textbf{R@5} & \textbf{R@10} & \textbf{mAP} & \textbf{mINP} \\ \hline
        TBPS-mSigLIP \cite{tran2025multilingual} & 46.73 & 68.65 & 77.55 & 41.75 & 26.56 \\ \hline
        Ours (Fixed Weight) & 56.87 & \textbf{77.18} & 84.15 & 50.70 & 34.61 \\ \hline
        \textbf{Ours (Curriculum)} & \textbf{57.10} & 76.98 & \textbf{84.34} & \textbf{50.90} & \textbf{34.85} \\ \hline
    \end{tabular}%
    }
\end{table}

The Curriculum strategy achieves the best performance on 4 out of 5 metrics, with only a marginal gap on R@5. This confirms that the warm-up phase prevents premature hard-negative mining from destabilizing alignment in low-data regimes.


\section{Conclusion and Future Work}\label{conclusion}
We propose a unified optimization framework for Multilingual Text-Based Person Search that combines Cross-modal Circle Loss with Low-Rank Adaptation (LoRA) to address the limitations of Sigmoid-based contrastive learning in low-resource settings. Circle Loss enhances embedding discrimination by alleviating gradient saturation for semi-hard negatives, while LoRA stabilizes optimization and enables larger batch sizes under constrained memory. Experiments across three typologically diverse languages (Vietnamese, English, Chinese) show that a curriculum hard-mining strategy achieves state-of-the-art results, including 52.28\% Rank-1 on 3000VnPersonSearch and 59.35\% on PRW-TPS-CN, while generalizing well to CUHK-PEDES. Future work includes exploring dynamic margin formulations to mitigate false negative collisions in high-resource datasets, extending to additional low-resource languages, and investigating alternative parameter-efficient fine-tuning strategies.

\section*{Acknowledgement}
This material is based upon work supported by the Air Force Office of Scientific Research under award number FA23862614005.

\vfill\pagebreak

\selectlanguage{English}
\bibliographystyle{IEEEbib}
\bibliography{Reference}

\end{document}
